% Notes of the Montreal seminar, 1970 (cote 129, typed part): archive cote 129, site page 35 (PDF page 36).
% typescript. Draft compared: none.
% AI-transcribed by Claude Opus 5.5 (claude-opus-5-5, adaptive thinking, effort high), Message Batches API, batch mtl129_full (msgbatch_01AGnDKHwy6ATCM7ehUa6263), 2026-11-13; not yet checked by eye.
\scanpage{129/35}
\mspage{\hand{II.7}}

2.1.3. \quad On dit que $G$ est \uline{unipotent} si $G^*$ est infinitésimal (voir
SGA 3, XVII, pour une définition générale des groupes unipotents).

2.1.4. \quad On dit que $G$ est de \uline{type multiplicatif} sur $S$ si $G^*$ est
étale. Cela équivaut au fait que pour tout $s \in S$ , $G_s$ est de type multi-%
plicatif, ou encore à ce que $G$ se plonge localement pour la topologie
étale dans $(\mathbb{G}_{m|S})^r$ .

\margin{X} 2.1.5. \quad On dit que $G$ est \uline{bi-infinitésimal} sur $S$ s'il est infini-%
tésimal et unipotent (i.e. $G$ et $G^*$ sont infinitésimaux).

2.2. \quad Si $G$ est étale et localement annulé par une puissance de $p$
sur un $F_p$-schéma $S$ , alors $G^*$ est infinitésimal. En d'autres termes,
sur une base de catactéristique $p$ , tout p-groupe de type multiplicatif
est infinitésimal. En effet, pour tout $s \in S$ , $G_{\overline{s}}$ est isomorphe au grou-%
pe constant $G(\overline{k})_{\overline{k}}$ (où l'on désigne par $\overline{k}$ la clôture algébrique du
corps résiduel en $s$). D'après 1.5.1 l'algèbre affine de $(G_{\overline{s}})^* = (G^*)_{\overline{s}}$
s'identifie à l'algèbre $\overline{k}[G(\overline{k})]$ du groupe $G(\overline{k})$ . On vérifie alors immé-%
diatement que l'idéal de $\overline{k}[G(\overline{k})]$ , formé des éléments $\sum\limits_{g\in G(\overline{k})} \lambda_g.g$ , $\lambda_g \in \overline{k}$
tels que $\sum\limits_{g} \lambda_g = 0$ , est maximal et nilpotent car $G(\overline{k})$ est annulé par une
puissance de $p$ et $\overline{k}$ est de caractéristique $p$ . Donc $(G^*)_{\overline{s}}$ se réduit
à un point et $G^*$ est infinitésimal.

2.3. \quad Soit
\[
0 \to G' \overset{\mu}{\to} G \overset{\pi}{\to} G'' \to 0
\]
une suite exacte de groupes finis localement libres. Alors $G$ est infini-%
tésimal (resp. étale, resp. unipotent, resp. de type multiplicatif, resp.

\margin{35}
